% Includable audit section; requires amsmath and amssymb.
\section{Independent audit of the direct progress rule}

Let $\mathcal C_D$ be the pinned native scalar calculus and let
$\mathcal C_Q$ have exactly the same terms, formulas, contexts, computation
axioms, transport, QE, DE, and case introduction. Replace only
\[
 \frac{\Gamma,E(t,z)\vdash E(t,Sz)}
      {\Gamma\vdash D_{x,y}E(x,y)}
 \quad(\mathrm{DI}_{\rm prog})
\]
by
\[
 \frac{\Gamma,E(z)\vdash E(Sz)}
      {\Gamma\vdash Q_yE(y)}
 \quad(\mathrm{QI}_{\rm prog}).
\]
Here $E$ is a literal equation. In either rule the induction eigenvariable
is fresh for the retained context and all fixed parameters; the old rule
also requires it to be absent from the witness $t$. The new rule contains
no existential witness or witness-substitution operation.

\paragraph{From the new rule to the old calculus.}
Choose a new variable $x$ absent from $E$. The old progress rule with
witness $0$ gives $D_{x,y}E(y)$ directly from the displayed premise.
Eliminate this formula at a fresh $b$, retaining $Q_yE(y)$, and return that
retained hypothesis. This is a three-node wrapper around the unchanged
step proof: old progress introduction, DE, and Hyp. Thus each new progress
node can be replaced by three old nodes.

\paragraph{From the old rule to the new calculus.}
Instantiate the old matrix first, writing $E_t(y)=E(t,y)$ with capture
avoiding substitution. The supplied step proof gives $Q_yE_t(y)$ by the
new rule. To introduce $D_{x,y}E(x,y)$ by the retained case rule, take
witness $t$, a new eigenvariable $w$, and local assumption $E_t(0)$.
QE applied to the derived $Q_yE_t(y)$ and that assumption gives
$E_t(Sw)$, as required. Weakening retains the local base assumption in the
derivation of $Q_yE_t(y)$. Its original induction eigenvariable remains
absent from the added base $E_t(0)$ by the old rule's freshness conditions:
it is absent from $t$ and all passive parameters of $E$. The new $w$ is
chosen outside the entire original proof. This wrapper consists of Q-progress,
QE, a hypothesis leaf, and case introduction. It replaces each old
progress node by four new nodes.

\paragraph{Full consequence equivalence.}
Both transformations preserve the concluding formula literally up to
alpha conversion and preserve every retained assumption. External
induction on a finite derivation therefore gives, for every native
context and formula,
\[
             \Gamma\vdash_{\mathcal C_D}A
             \quad\Longleftrightarrow\quad
             \Gamma\vdash_{\mathcal C_Q}A.
\]
This is consequence equivalence in the entire Eq/Q/D language, not just
agreement on equations or closed theorems. Proof translations are effective
finite tree transformations; they are not assertions of either theory's
internal consistency.

The ordinary numerical language and interpretation domain are unchanged.
Both calculi therefore present the same ordinary theory $T_{\rm eq}$.
In particular, the existing whole-domain bi-interpretation with iPRA
transfers with exactly the same term translations, formula translations,
actual equality, and identity graphs for both composites. Direct
soundness of Q-progress can also be checked in $T_{\rm eq}$: assume
$E(0)$, apply contextual equation induction to the supplied progress
derivation, universally generalize its conclusion, and discharge $E(0)$.
The result is precisely $E(0)\to\forall y\,E(y)$, the ordinary reading
of $Q_yE(y)$.

\paragraph{A nontrivial checked example.}
Put $F(n)=I(n,0;u.Su)$ and $E(y)\equiv F(y)=y$. The zero computation
gives $E(0)$. Under the induction assumption $F(z)=z$, successor
computation gives
\[
                       F(Sz)=S(F(z)).
\]
Transport in the template $F(Sz)=S(w)$ gives $F(Sz)=Sz$.
Q-progress followed by QE at the free parameter $n$ proves $F(n)=n$.
This displayed proof has six primitive nodes: a progress node, a transport
node, its equality hypothesis, its successor computation, a QE node, and
its zero-computation premise. The direct translation to the old calculus
has eight nodes. These are sizes of the specified proof trees, not claims
of shortest possible proofs.

\paragraph{Executable audit.}
The independent readable script \texttt{audit.py} checks all rules of both
calculi, performs capture avoiding substitution, compares bound syntax up
to alpha equivalence, and implements both progress-rule translations. It
checks the six-node example, the eight-node old translation, and the
eleven-node translation back. A separate nonvacuous-witness example uses
the equation $I(y,x;u.u)=x$, choosing a witness whose free name deliberately
matches the matrix's original index binder. The four-node old proof and
seven-node new proof both pass; the free witness parameter is preserved
by alpha renaming. Contextual translation and rejection of compound
progress, retained-context eigenvariable capture, and existential
eigenvariable escape are also checked. The retained-context fixture has
a correctly instantiated step; a rule-level ablation confirms that all
other Q-progress conditions pass when only freshness is omitted. The
production checker retains the freshness condition.

This script is supporting validation, not a trusted import of the compact
checker and not a proof-assistant formalization. The direct Q-progress
rule itself is already derived in the pinned Lean development as
\texttt{NPrf.qIProg}; the proposal is its selection as a primitive in place
of witness-bearing D-progress. No new arithmetic completeness theorem or
absolute minimality claim is inferred from that selection.
